\subsection{NEIGHBOURHOODS OF A POINT IN A TOPOLOGICAL GROUP}

Let \(\mathfrak{B}\) be the neighbourhood filter of the identity element \(e\) in a topological group G, and let \(a\) be any point of G. Since \(x \to ax\) and \(x \to xa\) are homeomorphisms, it follows that the neighbourhood filter of \(a\) is the family \(a\) . \(\mathfrak{B}\) of sets \(a\) . V, where V runs through \(\mathfrak{B}\) , and is also the family \(\mathfrak{B}\) . \(a\) of sets V. \(a\) . Thus we know the neighbourhood filter of any point of a topological group as soon as we know the neighbourhood filter of the identity element \(e\) of the group.

If we say that \(xy\) and \(x^{-1}\) are continuous at \(x = y = e\) , we obtain (Chapter I, § 2, no. 1):

\((\mathbf{GV}_{\mathbf{I}})\) . Given any \(\mathbf{U} \in \mathfrak{B}\) , there exists \(\mathbf{V} \in \mathfrak{B}\) such that \(\mathbf{V} \cdot \mathbf{V} \subset \mathbf{U}\) .

\((\mathbf{GV}_{\mathbf{II}})\) . Given any \(\mathbf{U}\in \mathfrak{B},\) we have \(\mathbf{U}^{-1}\in \mathfrak{B}\)

Every filter \(\mathfrak{B}\) on \(G\) which satisfies \((\mathrm{GV}_{\mathbb{I}})\) and \((\mathrm{GV}_{\mathbb{II}})\) also satisfies \((\mathrm{GV}_{\alpha})\) . Given any \(U \in \mathfrak{B}\) , there exists \(V \in \mathfrak{B}\) such that \(V \cdot V^{-1} \subset U\) . For by \((\mathrm{GV}_{\mathbb{I}})\) , there exists \(W \in \mathfrak{B}\) such that \(W \cdot W \subset U\) , and by \((\mathrm{GV}_{\mathbb{II}})\) there exists \(V \in \mathfrak{B}\) such that \(V \subset W \cap W^{-1}\) ; hence \(V^{-1} \subset W\) and therefore \(V \cdot V^{-1} \subset W \cdot W \subset U\) .

Conversely, if a filter \(\mathfrak{B}\) on \(\mathbf{G}\) satisfies \((\mathrm{GV}_{\alpha})\) , it follows first of all that \(e\) belongs to every set \(\mathbf{U} \in \mathfrak{B}\) ; for if \(\mathbf{V} \in \mathfrak{B}\) is such that \(\mathbf{V} \cdot \mathbf{V}^{-1} \subset \mathbf{U}\) , then, since V is not empty, we have \(x.x^{-1}=e\in U\) for every \(x\in V\) . The condition \((\mathrm{GV}_{\alpha})\) therefore implies that \(V^{-1}\subset V.V^{-1}\subset U\) , so that \(U^{-1}\in\mathfrak{B}\) whenever \(U\in\mathfrak{B}\) . Finally, if \(V\in\mathfrak{B}\) is such that \(V.V^{-1}\subset U\) , and \(W\in\mathfrak{B}\) is such that \(W\subset V\cap V^{-1}\) , we have \(W.W\subset U\) . We see thus that \((\mathrm{GV}_{\alpha})\) is equivalent to the conjunction of \((\mathrm{GV}_{I})\) and \((\mathrm{GV}_{II})\) .

Finally, since \(x \rightarrow axa^{-1}\) is a homeomorphism which leaves e fixed, B has the following property:

\((\mathrm{GV}_{\mathrm{III}})\) . For all \(a \in \mathbf{G}\) and all \(\mathbf{V} \in \mathfrak{B}\) , we have \(a. \mathbf{V}.a^{-1} \in \mathfrak{B}\) .

These three properties of the filter \(\mathfrak{B}\) are characteristic:

PROPOSITION 1. Let G be a group and let B be a filter on G satisfying the axioms \((\mathrm{GV}_{\mathrm{I}})\) , \((\mathrm{GV}_{\mathrm{II}})\) and \((\mathrm{GV}_{\mathrm{III}})\) . Then there is a unique topology on G, compatible with the group structure of G, for which B is the neighbourhood filter of the identity element e. For this topology the neighbourhood filter of any point \(a \in G\) is the same as each of the two filters a.B and B.a.

If there is a topology with the required properties, then by what has been said above the neighbourhood filter of \(a\) coincides with each of the filters \(a\) . \(\mathfrak{B}\) and \(\mathfrak{B}.a\) ; hence the topology is unique, if it exists. Its existence will be established if we show 1) that the filters \(a\) . \(\mathfrak{B}\) are the neighbourhood filters of a topology on \(\mathbf{G}\) , and 2) that this topology is compatible with the group structure of \(\mathbf{G}\) .

\begin{enumerate}
\def\labelenumi{\arabic{enumi})}
\item
  The filter \(a. \mathfrak{B}\) satisfies axiom \((\mathrm{V}_{\mathrm{III}})\) (see Chapter I § 1, no. 2) by reason of \((\mathrm{GV}_{\mathrm{I}})\) and \((\mathrm{GV}_{\mathrm{II}})\) , as we have already seen; hence to show that \(a. \mathfrak{B}\) is the neighbourhood filter of \(a\) in a topology on \(G\) , we have to verify axiom \((\mathrm{V}_{\mathrm{IV}})\) . Let then \(V\) be any set of \(\mathfrak{B}\) , and \(W\) a set of \(\mathfrak{B}\) such that \(W.W \subset V\) ; then for any \(x \in a.W\) we have \(x.W \subset a.W.W \subset a.V\) , so that \(a.V\) belongs to the filter \(x. \mathfrak{B}\) ; hence \((\mathrm{V}_{\mathrm{IV}})\) is satisfied.
\item
  Let us now show that the topology defined by the neighbourhood filters \(a\) . \(\mathfrak{B}\) satisfies (GT'). Let \(a, b\) be any two points of G; if we put \(x = au\) and \(y = bv\) , then we have to show that \(xy^{-1}\) is as near as we please to \(ab^{-1}\) whenever \(u\) and \(v\) are close enough to \(e\) . Now \((ab^{-1})^{-1}(xy^{-1}) = buv^{-1}b^{-1}\) ; let U be any neighbourhood of \(e\) , then we shall have \(buv^{-1}b^{-1} \in U\) if \(uv^{-1} \in b^{-1}Ub = V\) , and \(V \in \mathfrak{B}\) by reason of \((\mathrm{GV}_{\mathrm{III}})\) . But by \((\mathrm{GV}_{\mathrm{I}})\) and \((\mathrm{GV}_{\mathrm{II}})\) there exists \(W \in \mathfrak{B}\) such that \(W.W^{-1} \subset V\) ; hence it is enough to take \(u \in W\) and \(v \in W\) in order to have \(xy^{-1} \in (ab^{-1})U\) . This completes the proof.
\end{enumerate}

A common method of defining a topology compatible with a group structure on G consists in giving a filter satisfying the axioms \((\mathrm{GV}_{\mathrm{I}})\) , \((\mathrm{GV}_{\mathrm{II}})\) and \((\mathrm{GV}_{\mathrm{III}})\) . The corresponding conditions for a filter base B are as follows:

\((\mathbf{GV}_{\mathbf{I}}^{\prime})\) . Given any \(\mathbf{U}\in \mathfrak{B}\) , there exists \(\mathbf{V}\in \mathfrak{B}\) such that \(\mathbf{V}.\mathbf{V}\subset \mathbf{U}\) . \((\mathrm{GV}_{\mathrm{II}}^{\prime})\) . Given any \(\mathbf{U} \in \mathfrak{B}\) , there exists \(\mathbf{V} \in \mathfrak{B}\) such that \(\mathbf{V}^{-1} \subset \mathbf{U}\) . \((\mathrm{GV}_{\mathrm{III}}^{\prime})\) . Given any \(a \in \mathbf{G}\) and any \(\mathbf{U} \in \mathfrak{B}\) , there exists \(\mathbf{V} \in \mathfrak{B}\) such that \(\mathbf{V} \subset a. \mathbf{U}. a^{-1}\) .

A neighbourhood of e which coincides with its image under the symmetry \(x \rightarrow x^{-1}\) is said to be symmetric. If V is any neighbourhood of e, then \(V \cup V^{-1}\) , \(V \cap V^{-1}\) and \(V \cdot V^{-1}\) are symmetric neighbourhoods. By \((\mathrm{GV}_{\mathrm{II}})\) , the symmetric neighbourhoods form a fundamental system of neighbourhoods of e. Also it follows from \((\mathrm{GV}_{\mathrm{I}})\) that when V runs through a fundamental system of neighbourhoods of e, the sets \(V^{n}\) (where n is a fixed integer \(\neq 0\) ) form a fundamental system of neighbourhoods of e.

Remark. If G is commutative, we have \(x \cdot A \cdot x^{-1} = A\) for every subset A of G and every \(x \in G\) , and therefore \((GV_{III})\) {[}resp. \((GV_{III}')\) {]} is automatically satisfied for every filter (resp. filter base) on G. On the other hand, if G is not abelian, then \((GV_{III})\) is not a consequence of \((GV_I)\) and \((GV_{II})\) {[}see Exercise 5{]}.

If G is a commutative group, written additively, the axioms which characterize the filter B of neighbourhoods of the origin for a topology compatible with the group structure of G are therefore the following: \((\mathrm{GA}_{\mathrm{I}})\) . Given any \(U \in B\) , there exists \(V \in B\) such that \(V + V \subset U\) . \((\mathrm{GA}_{\mathrm{II}})\) . Given any \(U \in B\) , we have \(-U \in B\) .

PROPOSITION 2. A topological group G is Hausdorff if and only if the set \(\{e\}\) is closed.

Clearly, if G is Hausdorff, then \(\{e\}\) is closed. Conversely, if \(\{e\}\) is closed, then the diagonal \(\Delta\) of \(G \times G\) is closed, because it is the inverse image of \(\{e\}\) under the continuous mapping \((x, y) \to xy^{-1}\) ; hence (Chapter I, § 8, no. 1, Proposition 1) G is Hausdorff.

COROLLARY. A topological group G is Hausdorff if and only if the intersection of the neighbourhoods of e consists only of the point e.

The condition is clearly necessary. Conversely, if the intersection of all the neighbourhoods of \(e\) is just \(\{e\}\) , then given any \(x \neq e\) there is a neighbourhood \(V\) of \(e\) such that \(x^{-1} \notin V\) and therefore \(e \notin xV\) . This shows that \(x\) is not in the closure of \(\{e\}\) , and thus \(\{e\}\) is closed, so that \(G\) is Hausdorff.

Example. Definition of a topology on a group by means of a set of subgroups. If \(\mathfrak{B}\) is a filter base on a group \(\mathbf{G}\) , formed of subgroups of \(\mathbf{G}\) , then it is immediately seen that \(\mathfrak{B}\) satisfies axioms \((\mathrm{GV}_{\mathrm{I}}^{\prime})\) and \((\mathrm{GV}_{\mathrm{II}}^{\prime})\) , since \(\mathbf{H}.\mathbf{H}^{-1} = \mathbf{H}\) for any subgroup \(\mathbf{H}\) of \(\mathbf{G}\) . Hence the set \(\mathfrak{B}\) will be a fundamental system of neighbourhoods of \(e\) in a topology compatible with the group structure of \(\mathbf{G}\) , provided that \(\mathfrak{B}\) satisfies \((\mathrm{GV}_{\mathrm{III}})\) ; this will in particular be the case if all the subgroups in \(\mathfrak{B}\) are normal, hence always if \(G\) is commutative. The topology thus defined is Hausdorff, by Proposition 2, if and only if the intersection of all the subgroups in \(\mathfrak{B}\) consists only of \(e\) . The most interesting cases are those in which the subgroup \(\{e\}\) is not in \(\mathfrak{B}\) (otherwise the topology defined by \(\mathfrak{B}\) is the discrete topology): if \(\{e\} \notin \mathfrak{B}\) , the topology defined by \(\mathfrak{B}\) is Hausdorff only if \(\mathfrak{B}\) is an infinite set.

Since the intersection of two subgroups is a subgroup, we can define a topology on G, compatible with its group structure, starting from any set \(\mathfrak{F}\) of subgroups of G: let \(\mathfrak{G}\) be the set of all subgroups \(a.H.a^{-1}\) , where \(H\in\mathfrak{F}\) and \(a\in G\) , and let \(\mathfrak{B}\) be the set of all finite intersections of subgroups belonging to \(\mathfrak{G}\) . Then \(\mathfrak{B}\) is a filter base and satisfies \((\mathrm{GV}_{\mathrm{III}}^{\prime})\) .

Consider in particular the additive group of a ring A. Every set \(\mathfrak{F}\) of ideals of A defines a topology compatible with this additive group structure. This topology is Hausdorff if the intersection of all the ideals of \(\mathfrak{F}\) is the zero ideal, and it is not discrete if no finite intersection of the ideals of \(\mathfrak{F}\) is the zero ideal. Topologies defined in this way play a large part in the theory of numbers (see the Exercises of §§ 6 and 7 of this chapter).

